\begin{definition}[Linearly independent operator family]%
    \label{def:mpdo_bnt_label_linear_independent}
    \lean{MPOTensor.BNTLabelOperatorFamily.LinearIndependentAt}
    \lean{MPOTensor.BNTLabelOperatorFamily.EventuallyLinearIndependent}
    \leanok
    \uses{def:mpdo_bnt_label_operator_family}
    The family $\{O_L(M_\alpha)\}_\alpha$ is \emph{linearly independent at length $L$}
    when its operators are linearly independent, and
    \emph{eventually linearly independent} when this holds for every length
    beyond some threshold. Eventual linear independence is the
    operator form of the third defining property of a basis of normal
    tensors in \cite[Section~2.3]{Cirac2016MPDO_arXiv}, invoked for the
    vertical operators in \cite[Appendix~C.4]{Cirac2016MPDO_arXiv}.
\end{definition}

\begin{theorem}[Uniqueness of the structure coefficients]%
    \label{thm:mpdo_bnt_label_coeff_unique}
    \lean{MPOTensor.BNTLabelOperatorFamily.HasSameLengthProductForm.%
        coeff_eq_of_linearIndependentAt}
    \leanok
    \uses{def:mpdo_bnt_label_linear_independent, def:mpdo_bnt_label_same_length_product_form, def:mpdo_bnt_label_coefficient_family}
    Two systems of structure coefficients satisfying the product law for the
    same operator family agree at every positive length at which the
    operators are linearly independent.
\end{theorem}
\begin{proof}\leanok
    Applying the product law to both systems of structure coefficients gives
    \begin{align}
        \sum_\delta c^{(L)}_{\alpha,\beta,\delta}\,O_L(M_\delta)
         & =
        O_L(M_\alpha)O_L(M_\beta)
        =
        \sum_\delta c'^{(L)}_{\alpha,\beta,\delta}\,O_L(M_\delta).
        \notag
    \end{align}
    Subtracting yields
    $\sum_\delta (c^{(L)}_{\alpha,\beta,\delta}
        - c'^{(L)}_{\alpha,\beta,\delta})O_L(M_\delta) = 0$,
    and linear independence forces every coefficient of the difference to
    vanish.
\end{proof}

\begin{definition}[Change of vertical representatives]%
    \label{def:mpdo_vertical_bnt_explicit_transport}
    \lean{MPOTensor.BNTLabelOperatorFamily.ExplicitVerticalTransport}
    \leanok
    \uses{def:mpdo_bnt_label_operator_family}
    Let $\{O_L(M_\alpha)\}_{\alpha\in\Lambda}$ and
    $\{O_L(M'_\alpha)\}_{\alpha\in\Lambda'}$ be two operator families.
    A change of vertical representatives consists of a bijection
    $e:\Lambda'\to\Lambda$, nonzero complex scalars $s_\alpha$, and complex algebra
    isomorphisms $\Phi_L$ between the ambient operator algebras.
    For every positive length $L$ and every $\alpha\in\Lambda'$, these satisfy
    \begin{align}
        O_L(M'_\alpha) & =s_\alpha^L\Phi_L(O_L(M_{e(\alpha)})).
        \notag
    \end{align}
\end{definition}

\begin{definition}[Structure coefficients under a change of normalization]%
    \label{def:mpdo_vertical_bnt_transported_coefficients}
    \lean{MPOTensor.BNTLabelOperatorFamily.verticalTransportCoefficients}
    \leanok
    \uses{def:mpdo_bnt_label_coefficient_family}
    Given a label bijection $e:\Lambda'\to\Lambda$ and complex scalars $s_\alpha$,
    define
    \begin{align}
        \widetilde c^{(L)}_{\alpha,\beta,\gamma}
         & =\left(\frac{s_\alpha s_\beta}{s_\gamma}\right)^L
        c^{(L)}_{e(\alpha),e(\beta),e(\gamma)}.
        \notag
    \end{align}
    For this definition, use the convention $z/0=0$.
\end{definition}

\begin{theorem}[Product law under a change of vertical representatives]%
    \label{thm:mpdo_vertical_bnt_product_transport}
    \lean{MPOTensor.BNTLabelOperatorFamily.ExplicitVerticalTransport.%
        hasSameLengthProductForm_verticalTransport}
    \leanok
    \uses{def:mpdo_vertical_bnt_explicit_transport, def:mpdo_vertical_bnt_transported_coefficients, def:mpdo_bnt_label_same_length_product_form}
    A change of vertical representatives carries a product law with
    structure coefficients $c^{(L)}_{\alpha,\beta,\gamma}$ to a product law
    with structure coefficients $\widetilde c^{(L)}_{\alpha,\beta,\gamma}$.
\end{theorem}
\begin{proof}\leanok
    For every positive length $L$, multiplication of the two new
    representatives and the original product law give
    \begin{align}
        O_L(M'_\alpha)O_L(M'_\beta)
         & =(s_\alpha s_\beta)^L
        \Phi_L(O_L(M_{e(\alpha)})O_L(M_{e(\beta)}))
        \notag                   \\
         & =(s_\alpha s_\beta)^L
        \sum_\delta c^{(L)}_{e(\alpha),e(\beta),\delta}
        \Phi_L(O_L(M_\delta)).
        \notag
    \end{align}
    Reindexing by $\delta=e(\gamma)$ and using $s_\gamma\ne0$ yields
    \begin{align}
        O_L(M'_\alpha)O_L(M'_\beta)
         & =\sum_\gamma
        \left(\frac{s_\alpha s_\beta}{s_\gamma}\right)^L
        c^{(L)}_{e(\alpha),e(\beta),e(\gamma)}
        \,s_\gamma^L\Phi_L(O_L(M_{e(\gamma)}))
        \notag          \\
         & =\sum_\gamma
        \widetilde c^{(L)}_{\alpha,\beta,\gamma}O_L(M'_\gamma).
        \notag
    \end{align}
\end{proof}

\begin{theorem}[Covariance of structure coefficients for operator families]%
    \label{thm:mpdo_vertical_bnt_coefficient_covariance_abstract}
    \lean{MPOTensor.BNTLabelOperatorFamily.ExplicitVerticalTransport.%
        coeff_eq_verticalTransport_of_linearIndependentAt}
    \leanok
    \uses{def:mpdo_bnt_label_linear_independent, thm:mpdo_bnt_label_coeff_unique, thm:mpdo_vertical_bnt_product_transport}
    Suppose that two operator families satisfy product laws and are related
    by a change of vertical representatives. At a positive length where
    the operators of the second family are linearly independent, their
    structure coefficients equal the transformed coefficients
    $\widetilde c^{(L)}_{\alpha,\beta,\gamma}$ of the first family.
\end{theorem}
\begin{proof}\leanok
    The product law for the second family and
    Theorem~\ref{thm:mpdo_vertical_bnt_product_transport} give
    \begin{align}
        \sum_\gamma c'^{(L)}_{\alpha,\beta,\gamma}O_L(M'_\gamma)
         & =O_L(M'_\alpha)O_L(M'_\beta)
        =\sum_\gamma
        \widetilde c^{(L)}_{\alpha,\beta,\gamma}O_L(M'_\gamma).
        \notag
    \end{align}
    Linear independence of these operators identifies the two
    systems of structure coefficients term by term.
\end{proof}

\begin{theorem}[Covariance of structure coefficients for vertical BNT decompositions]%
    \label{thm:mpdo_vertical_bnt_coefficient_covariance}
    \lean{MPOTensor.BNTAlgebraTensorClause.%
        coeff_eq_of_explicitVerticalTransport}
    \leanok
    \uses{def:mpdo_bnt_algebra_tensor_clause, def:mpdo_vertical_bnt_explicit_transport, def:mpdo_bnt_label_linear_independent, thm:mpdo_vertical_bnt_coefficient_covariance_abstract}
    Suppose that two vertical BNT decompositions of tensors satisfy the
    algebra condition and are related by a change of vertical representatives.
    At every positive length $L$ at which the operators of the second
    decomposition are linearly independent, their structure coefficients satisfy
    \begin{align}
        {c'}^{(L)}_{\alpha,\beta,\gamma}
         & =\left(\frac{s_\alpha s_\beta}{s_\gamma}\right)^L
        c^{(L)}_{e(\alpha),e(\beta),e(\gamma)}.
        \notag
    \end{align}
    Theorem~\ref{thm:mpdo_vertical_coefficients_not_presentation_invariant}
    shows why equality of the horizontal periodic operator families alone
    cannot replace the specified relation between vertical representatives.
\end{theorem}
\begin{proof}\leanok
    The algebra conditions for the two vertical BNT decompositions give the
    product laws required by the covariance theorem for operator families.
    Hence
    \begin{align}
        {c'}^{(L)}_{\alpha,\beta,\gamma}
         & =\widetilde c^{(L)}_{\alpha,\beta,\gamma}
        =\left(\frac{s_\alpha s_\beta}{s_\gamma}\right)^L
        c^{(L)}_{e(\alpha),e(\beta),e(\gamma)}.
        \notag
    \end{align}
\end{proof}

\begin{theorem}[Equality of structure coefficients under compatible normalization]%
    \label{thm:mpdo_vertical_bnt_coefficient_transport_exact_abstract}
    \lean{MPOTensor.BNTLabelOperatorFamily.ExplicitVerticalTransport.%
        coeff_eq_of_scale_mul_eq_of_linearIndependentAt}
    \leanok
    \uses{thm:mpdo_vertical_bnt_coefficient_covariance_abstract}
    Under the hypotheses of
    Theorem~\ref{thm:mpdo_vertical_bnt_coefficient_covariance_abstract}, if
    $s_\alpha s_\beta=s_\gamma$, then
    ${c'}^{(L)}_{\alpha,\beta,\gamma}
        =c^{(L)}_{e(\alpha),e(\beta),e(\gamma)}$.
\end{theorem}
\begin{proof}\leanok
    The compatible normalization and $s_\gamma\ne0$ give
    $\left(s_\alpha s_\beta/s_\gamma\right)^L
        =\left(s_\gamma/s_\gamma\right)^L=1$.
    Substitution in the covariance formula gives the result.
\end{proof}

\begin{theorem}[Equality of vertical BNT structure coefficients under compatible normalization]%
    \label{thm:mpdo_vertical_bnt_coefficient_transport_exact}
    \lean{MPOTensor.BNTAlgebraTensorClause.%
        coeff_eq_of_explicitVerticalTransport_of_scale_mul_eq}
    \leanok
    \uses{def:mpdo_bnt_algebra_tensor_clause, thm:mpdo_vertical_bnt_coefficient_transport_exact_abstract}
    Under the hypotheses of
    Theorem~\ref{thm:mpdo_vertical_bnt_coefficient_covariance}, if
    $s_\alpha s_\beta=s_\gamma$, then
    ${c'}^{(L)}_{\alpha,\beta,\gamma}
        =c^{(L)}_{e(\alpha),e(\beta),e(\gamma)}$.
\end{theorem}
\begin{proof}\leanok
    For the vertical BNT decompositions, the covariance formula gives
    \begin{align}
        {c'}^{(L)}_{\alpha,\beta,\gamma}
         & =\left(\frac{s_\alpha s_\beta}{s_\gamma}\right)^L
        c^{(L)}_{e(\alpha),e(\beta),e(\gamma)}
        =c^{(L)}_{e(\alpha),e(\beta),e(\gamma)},
        \notag
    \end{align}
    where the last equality uses $s_\alpha s_\beta=s_\gamma$ and
    $s_\gamma\ne0$.
\end{proof}

% Project-derived counterexample to the cross-presentation assertion formerly
% proposed in #6395. CPSV16 does not state this assertion.
\begin{definition}[Horizontally gauge-equivalent tensors with distinct vertical normalizations]%
    \label{def:mpdo_vertical_coefficient_presentation_counterexample}
    \lean{MPOTensor.VerticalCoefficientPresentationCounterexample.leftMPO}
    \lean{MPOTensor.VerticalCoefficientPresentationCounterexample.rightMPO}
    \lean{MPOTensor.VerticalCoefficientPresentationCounterexample.leftClause}
    \lean{MPOTensor.VerticalCoefficientPresentationCounterexample.rightClause}
    \leanok
    \uses{def:mpo_tensor, def:mpdo_bnt_algebra_tensor_clause}
    Let
    \begin{align}
        P & =\begin{pmatrix}1&0\\0&0\end{pmatrix},
        \notag                                        \\
        X & =\begin{pmatrix}1&-3/4\\0&1\end{pmatrix},
        \notag                                        \\
        Q & =XPX^{-1}
        =\begin{pmatrix}1&3/4\\0&0\end{pmatrix}.
        \notag
    \end{align}
    Define MPO tensors with one-dimensional physical space by
    $M_P^{00}=P$ and $M_Q^{00}=Q$. Equip them with single-element vertical
    BNT decompositions whose normalized representatives and positive
    multiplicity weights are
    \begin{align}
        A_P & =P,
        \notag           \\
        m_P & =1,
        \notag           \\
        A_Q & =\frac45Q,
        \notag           \\
        m_Q & =\frac54.
        \notag
    \end{align}
    The corresponding positive diagonal matrices $\chi$ in the algebra condition are
    $[1]$ and $[4/5]$. No horizontal canonical-form assumption is imposed
    on the sheared tensor $M_Q$.
\end{definition}

\begin{theorem}[Horizontal periodic equality alone does not determine vertical BNT coefficients]%
    \label{thm:mpdo_vertical_coefficients_not_presentation_invariant}
    \lean{MPOTensor.VerticalCoefficientPresentationCounterexample.%
        leftMPO_gaugeEquiv_rightMPO}
    \lean{MPOTensor.VerticalCoefficientPresentationCounterexample.%
        leftMPO_sameMPV₂Pos_rightMPO}
    \lean{MPOTensor.VerticalCoefficientPresentationCounterexample.leftMPO_isMPDO}
    \lean{MPOTensor.VerticalCoefficientPresentationCounterexample.rightMPO_isMPDO}
    \lean{MPOTensor.VerticalCoefficientPresentationCounterexample.leftClause_coeff}
    \lean{MPOTensor.VerticalCoefficientPresentationCounterexample.rightClause_coeff}
    \lean{MPOTensor.VerticalCoefficientPresentationCounterexample.%
        sameMPV₂Pos_and_no_relabelled_coefficient_equality}
    \leanok
    \uses{def:mpdo, def:gauge_equiv, def:same_mpv2_pos, def:mpdo_vertical_coefficient_presentation_counterexample}
    The tensors $M_P$ and $M_Q$ are related by the nonunitary gauge $X$.
    For every positive length $N$, their unnormalized periodic operators
    satisfy $\rho^{(N)}(M_P)=\rho^{(N)}(M_Q)$.
    In particular, both tensors generate matrix product density operators.
    Their chosen vertical BNT decompositions nevertheless have structure
    coefficients
    \begin{align}
        c_P^{(L)} & =1,
        \notag                               \\
        c_Q^{(L)} & =\left(\frac45\right)^L.
        \notag
    \end{align}
    Hence there is no relabelling of the singleton label sets for which,
    at every positive length,
    \begin{align}
        c_Q^{(L)}(\alpha,\beta,\gamma)
         & =c_P^{(L)}(e(\alpha),e(\beta),e(\gamma)).
        \notag
    \end{align}
    Indeed, the two sides at length one are $4/5$ and $1$.

    Proposition~4.13 and Theorem~4.14 of
    \cite{Cirac2016MPDO_arXiv} instead start with one horizontally canonical
    tensor generating MPDOs and choose its vertical decomposition.
    The tensor $M_Q$ is not asserted to be a horizontal canonical
    representative. Thus this counterexample excludes a conclusion about
    distinct vertical decompositions based on periodic equality alone; it
    does not alter the source's characterization for a fixed decomposition.
    It also leaves
    Theorem~\ref{thm:mpdo_bnt_label_coeff_unique} unchanged, since that theorem
    compares two systems of structure coefficients for the same vertical
    operator family.
\end{theorem}
\begin{proof}\leanok
    The identities $P^2=P$, $Q^2=Q$, and
    $\tr(P)=\tr(Q)=1$ prove the periodic equality. The Frobenius norms of
    $P$ and $Q$ are $1$ and $5/4$, respectively, which gives the
    normalized representatives and multiplicity weights. Their products
    satisfy
    \begin{align}
        A_P^2 & =A_P,
        \notag               \\
        A_Q^2 & =\frac45A_Q.
        \notag
    \end{align}
    Taking the $L$-fold vertical operators gives
    $c_P^{(L)}=1$ and $c_Q^{(L)}=(4/5)^L$.
    Their values at $L=1$ exclude the only possible relabelling.
\end{proof}

\begin{theorem}[Length independence is not preserved by horizontal periodic equality alone]%
    \label{thm:mpdo_vertical_length_independence_not_presentation_invariant}
    \lean{MPOTensor.VerticalCoefficientPresentationCounterexample.%
        leftClause_lengthIndependent}
    \lean{MPOTensor.VerticalCoefficientPresentationCounterexample.%
        rightClause_not_lengthIndependent}
    \lean{MPOTensor.VerticalCoefficientPresentationCounterexample.%
        sameMPV₂Pos_and_lengthIndependent_mismatch}
    \leanok
    \uses{def:same_mpv2_pos, def:mpdo_bnt_label_length_independent, def:mpdo_vertical_coefficient_presentation_counterexample}
    The structure coefficients in the vertical decomposition of $M_P$ are
    length independent, whereas those in the chosen vertical decomposition
    of $M_Q$ are not. More precisely,
    \begin{align}
        c_P^{(L)} & =1,
        \notag                      \\
        c_Q^{(1)} & =\frac45,
        \notag                      \\
        c_Q^{(2)} & =\frac{16}{25}.
        \notag
    \end{align}
\end{theorem}
\begin{proof}\leanok
    The first family is constant, whereas $c_Q^{(1)}\ne c_Q^{(2)}$.
\end{proof}

\begin{theorem}[Associativity constraint on the structure coefficients]%
    \label{thm:mpdo_bnt_label_coeff_associativity}
    \lean{MPOTensor.BNTLabelOperatorFamily.HasSameLengthProductForm.%
        coeff_assoc_of_linearIndependentAt}
    \leanok
    \uses{def:mpdo_bnt_label_linear_independent, def:mpdo_bnt_label_same_length_product_form}
    At every positive length $L$ at which the operators are linearly
    independent, the structure coefficients of the product law
    satisfy
    \begin{align}
        \sum_\delta
        c^{(L)}_{\alpha,\beta,\delta}
        c^{(L)}_{\delta,\gamma,\epsilon}
         & =
        \sum_\delta
        c^{(L)}_{\beta,\gamma,\delta}
        c^{(L)}_{\alpha,\delta,\epsilon}.
        \notag
    \end{align}
    These are the restrictions that associativity of multiplication
    imposes on the structure coefficients, noted in
    \cite[Section~4.5]{Cirac2016MPDO_arXiv}.
\end{theorem}
\begin{proof}\leanok
    Applying the product law twice in each order gives
    \begin{align}
        (O_L(M_\alpha)O_L(M_\beta))O_L(M_\gamma)
         & =
        \sum_{\varepsilon'}
        \left(\sum_\delta
        c^{(L)}_{\alpha,\beta,\delta}
        c^{(L)}_{\delta,\gamma,\varepsilon'}\right)
        O_L(M_{\varepsilon'}),
        \notag \\
        O_L(M_\alpha)(O_L(M_\beta)O_L(M_\gamma))
         & =
        \sum_{\varepsilon'}
        \left(\sum_\delta
        c^{(L)}_{\beta,\gamma,\delta}
        c^{(L)}_{\alpha,\delta,\varepsilon'}\right)
        O_L(M_{\varepsilon'}).
        \notag
    \end{align}
    Associativity of multiplication equates the two left-hand sides, so the
    right-hand sides agree; linear independence of the operators then forces
    the corresponding coefficient sums to agree for each $\varepsilon'$.
\end{proof}

\begin{theorem}[Associativity constraint in trace-power form]%
    \label{thm:mpdo_bnt_label_chi_associativity}
    \lean{MPOTensor.BNTLabelOperatorFamily.HasSameLengthProductForm.%
        tracePowerCoeff_assoc_of_linearIndependentAt}
    \leanok
    \uses{thm:mpdo_bnt_label_coeff_associativity, def:mpdo_bnt_label_same_length_product_form, def:mpdo_bnt_label_positive_length_chi_trace_power_form}
    Suppose that the operators $\{O_L(M_\alpha)\}_\alpha$ satisfy the product
    law with structure coefficients
    $c^{(L)}_{\alpha,\beta,\gamma}=\tr(\chi_{\alpha,\beta,\gamma}^{\,L})$
    for diagonal matrices $\chi_{\alpha,\beta,\gamma}$.
    At every positive length $L$ at which the operators are linearly
    independent, the associativity constraint becomes
    \begin{align}
        \sum_\delta
        \tr(\chi_{\alpha,\beta,\delta}^{\,L})
        \tr(\chi_{\delta,\gamma,\epsilon}^{\,L})
         & =
        \sum_\delta
        \tr(\chi_{\beta,\gamma,\delta}^{\,L})
        \tr(\chi_{\alpha,\delta,\epsilon}^{\,L}).
        \notag
    \end{align}
    In \cite[Section~4.5]{Cirac2016MPDO_arXiv}, these restrictions in the
    positive diagonal setting, together with the isometry statement, lead
    to a pentagon-like equation for the fusion isometries. Solutions
    constructed from unitary fusion categories give the family of RFP MPDOs
    described there.
\end{theorem}
\begin{proof}\leanok
    Substitute the trace-power form of the coefficients into the
    associativity constraint.
\end{proof}

\begin{theorem}[The dimer RFP has length-dependent coefficients in its chosen vertical BNT decomposition]%
    \label{thm:mpdo_bnt_label_rescaling_stable_rfp_capstone}
    \lean{MPOTensor.RescalingStableLengthDependentRFP.%
        exists_isRFPViaTS_not_lengthIndependent_bntCoefficients_R}
    \leanok
    \uses{def:cpsv_canonical_form, def:mpdo, def:is_rfp_via_ts, def:mpdo_bnt_algebra_tensor_clause, def:mpdo_bnt_label_length_independent}
    There is a tensor $M$ generating MPDOs that has the direct-sum form
    $M^i=\bigoplus_k\mu_k M_k^i$ with normalized NT $M_k$, is a
    renormalization fixed point, and has a vertical BNT decomposition
    satisfying the algebra condition whose structure coefficients depend on
    the positive length $L$. One may take $M=R$ and the single-element BNT
    decomposition already constructed, for which
    \begin{align}
        c^{(L)}_{0,0,0} & =1+\left(\frac{7}{25}\right)^L.
        \notag
    \end{align}
    This exhibits length-dependent structure coefficients for the chosen
    vertical decomposition of $R$, addressing the question following
    Theorem~4.14 of \cite{Cirac2016MPDO_arXiv} without requiring unimodular
    canonical weights.
    % Scope note: docs/paper-gaps/cpsv16_unit_weight_rfp_scale_tension.tex
\end{theorem}

\begin{proof}\leanok
    \uses{thm:mpdo_bnt_label_rescaling_stable_length_dependence, thm:mpdo_bnt_label_rescaling_stable_cf, thm:mpdo_bnt_label_rescaling_stable_ismpdo, thm:mpdo_bnt_label_rescaling_stable_is_rfp_via_ts, thm:mpdo_bnt_label_rescaling_stable_bnt_algebra_clause}
    The preceding results show that $R$ generates MPDOs and has the
    direct-sum form with normalized NT blocks. They construct the tpCPM
    $S,T$ satisfying $S[R_2(X)]=R_1(X)$ and $T[R_1(X)]=R_2(X)$ for all
    $X$, together with its single-element vertical BNT decomposition
    satisfying the algebra condition. Its structure coefficients are
    $c^{(L)}_{0,0,0}=1+(7/25)^L$. The values at lengths one and two differ,
    so these coefficients are not length independent.
\end{proof}
